% ============================================================
% preamble.tex — configuración compartida (sin \documentclass)
% ============================================================

% ---------- Idioma y fuentes ----------
\usepackage[T1]{fontenc}
\usepackage{lmodern}
\usepackage[spanish,es-noquoting,es-nodecimaldot,es-tabla,es-lcroman]{babel}
\usepackage[autostyle]{csquotes}
\usepackage{microtype}

% ---------- Matemática ----------
\usepackage{amsmath}
\usepackage{amssymb}
\usepackage{siunitx}
\sisetup{per-mode=symbol,range-units=single}

% ---------- Colores ----------
\usepackage{xcolor}
\definecolor{tecnavy}{RGB}{31,56,100}
\definecolor{tecsteel}{RGB}{46,85,150}
\definecolor{teclight}{RGB}{214,227,243}
\definecolor{tecrule}{RGB}{155,184,216}
\definecolor{boxbg}{RGB}{242,246,252}
\definecolor{labelgray}{RGB}{120,120,120}
\definecolor{notaverde}{RGB}{22,110,60}
\definecolor{decisionambar}{RGB}{170,110,10}

% ---------- Gráficos ----------
\usepackage{graphicx}
\graphicspath{{img/}{figures/}{./}}
\usepackage{tikz}
\usetikzlibrary{calc,babel}

% ---------- Tablas y listas ----------
\usepackage{float}
\usepackage{array}
\usepackage{booktabs}
\usepackage{tabularx}
\usepackage[justification=centering]{caption}
\usepackage{enumitem}
\newcolumntype{P}[1]{>{\raggedright\arraybackslash}p{#1}}
\newcolumntype{Y}{>{\raggedright\arraybackslash}X}
\renewcommand{\arraystretch}{1.15}
\setlist[itemize]{leftmargin=*,noitemsep,topsep=2pt}
\setlist[enumerate]{leftmargin=*,noitemsep,topsep=2pt}

% ---------- Página ----------
\usepackage[a4paper,margin=1in,headheight=15pt]{geometry}
\usepackage{setspace}
\onehalfspacing
\usepackage{ragged2e}
\usepackage{fancyhdr}
\usepackage{lastpage}
\pagestyle{fancy}
\fancyhf{}
\fancyhead[L]{\footnotesize\textsc{\docCode}}
\fancyhead[R]{\footnotesize\docShortTitle}
\fancyfoot[C]{\footnotesize P\'agina \thepage\ de \pageref{LastPage}}
\renewcommand{\headrulewidth}{0.4pt}
\renewcommand{\footrulewidth}{0pt}

% ---------- Enlaces ----------
\usepackage{hyperref}
\hypersetup{
  colorlinks=true,
  linkcolor=tecnavy,
  urlcolor=tecsteel,
  pdftitle={Atributo de Investigaci\'on (IN) - Tema 9},
  pdfauthor={\docAuthors},
  pdfsubject={\docCourse}
}

% ---------- Entrada segura ----------
\newcommand{\cajaarchivofaltante}[1]{%
  \par\medskip\noindent
  \fcolorbox{red!60!black}{red!5}{%
    \parbox{\dimexpr\linewidth-2\fboxsep-2\fboxrule}{%
      \centering\ttfamily\small ARCHIVO FALTANTE: \detokenize{#1}}}%
  \par\medskip}

\newcommand{\cajaimagenfaltante}[1]{%
  \begingroup\setlength{\fboxsep}{0pt}%
  \fcolorbox{tecrule}{boxbg}{%
    \parbox[c][2.8cm][c]{6cm}{\centering\ttfamily\scriptsize
      IMAGEN NO ENCONTRADA\\[3pt]\detokenize{#1}}}%
  \endgroup}

\newcommand{\safeinput}[1]{%
  \IfFileExists{#1.tex}{\input{#1}}{%
    \PackageWarning{tecdoc}{Falta el archivo de secci\'on #1.tex}%
    \cajaarchivofaltante{#1.tex}}}

\newcommand{\safeimage}[2][width=0.75\linewidth]{%
  \IfFileExists{#2}{\includegraphics[#1]{#2}}{%
    \IfFileExists{img/#2}{\includegraphics[#1]{img/#2}}{%
      \PackageWarning{tecdoc}{Falta la imagen #2}%
      \cajaimagenfaltante{#2}}}}

\newcommand{\code}[1]{\texttt{#1}}

% ---------- Bloques del atributo IN ----------
% \indicadorIN{código}{texto literal del indicador}
% La respuesta del grupo se escribe inmediatamente después.
\newcommand{\indicadorIN}[2]{%
  \subsection{Indicador #1}
  \noindent\fcolorbox{tecrule}{boxbg}{%
    \parbox{\dimexpr\linewidth-2\fboxsep-2\fboxrule}{%
      \vspace{3pt}\textbf{\color{tecnavy}Texto del indicador.}\ #2\vspace{3pt}}}%
  \par\medskip\noindent
  \textbf{\color{tecnavy}Respuesta del grupo.}\ }

% ---------- Marcadores de redacción ----------
\newcommand{\pc}[1]{%
  \texorpdfstring{\textcolor{red}{\textbf{[POR COMPLETAR: #1]}}}%
                 {[POR COMPLETAR: #1]}}
\newcommand{\razon}[1]{%
  \texorpdfstring{\textcolor{notaverde}{\small\textit{[RAZ\'ON: #1]}}}%
                 {[RAZON: #1]}}
\newcommand{\decidir}[1]{%
  \texorpdfstring{\textcolor{decisionambar}{\small\textbf{[DECIDIR: #1]}}}%
                 {[DECIDIR: #1]}}

% ---------- Apagado de marcadores (descomentar antes de entregar) ----------
%\renewcommand{\pc}[1]{}
%\renewcommand{\razon}[1]{}
%\renewcommand{\decidir}[1]{}